\subsection{Causal interventions on heads}
\label{sec:rw-patching}

\emph{Activation patching} replaces the activation of one component on a clean input
with its activation on a minimally different corrupted input and measures the change
in the model's output \citep{wang2023ioi,heimersheim2024patching}. Run in the
\emph{noising} direction (corrupted activation into the clean run) it measures how
much of the output depends on the component; run in the \emph{denoising} direction
(clean activation into the corrupted run) it measures how much the component alone
restores. \citet{heimersheim2024patching} note that the two are not symmetric and
recommend reporting both. The choice of metric matters as much as the intervention:
\citet{zhang2024bestpractices} show that logit difference between the correct and the
corrupted answer is more informative and less brittle than probability or accuracy,
and that conclusions about which components matter can flip between metrics. The
indirect-object-identification (IOI) study of \citet{wang2023ioi} established the
per-head, per-position form of these tools on GPT-2 small and introduced the
vocabulary that later work inherited: \emph{name movers} that copy the answer,
\emph{negative name movers} that suppress it, and \emph{backup} heads that activate
only when others are removed. We use noising and denoising head patching with the
logit-difference metric throughout (\S\ref{sec:setting}), at all prompt positions and
at the answer position separately, and we split a head's patched output into the
attention pattern and the values it reads, as in the QK/OV decomposition of
\S\ref{sec:sih}.

\subsection{Reading what a head writes}
\label{sec:rw-readout}

Intervention tells us that a head matters; a second family of tools asks what it
contributes. \emph{Direct logit attribution} projects a head's output vector onto the
unembedding to see which tokens it promotes \citep{elhage2021mathematical,wang2023ioi};
the relation index of \S\ref{sec:sih} is a weight-only version of this projection that
uses the current token's embedding instead of the head's actual output. Behavioural
\emph{fingerprints} test a head on controlled synthetic inputs: prefix-matching and
copying scores on repeated random sequences \citep{olsson2022induction} and
needle-retrieval scores on key--value contexts \citep{wu2024retrieval} identify
induction and retrieval heads independently of any natural-language task. Finally,
\emph{Patchscopes} \citep{ghandeharioun2024patchscopes} decode a hidden
representation by injecting it into a separate prompt whose continuation exposes its
content; the method reads representations, not head contributions, so head
attribution requires comparing intact and head-intervened representations under the
same readout. We use all three: the contextual output projection and its
raw-embedding counterpart side by side, both synthetic fingerprints on every head in
our inventory, and a controlled Patchscopes-style readout of the sites where our
strongest head writes.

\subsection{From heads to circuits}
\label{sec:rw-circuits}

A list of important heads is not a mechanism. \emph{Path patching}
\citep{wang2023ioi} isolates a single edge: the output of a source component is
patched only along its path into a chosen receiver (for a head, its query, key or
value input), so an effect on the output shows that the receiver uses what the source
wrote. \citet{conmy2023acdc} automate this search, pruning edges whose removal does
not change a task metric, and \citet{haklay2025position} show that edges must be
indexed by token position: the same head can be essential at one position and inert
at another, and position-agnostic circuits miss this. How a discovered circuit is
judged also matters. \citet{hanna2024faithfulness} argue that overlap with a
hand-found reference is the wrong target and that a circuit should be evaluated by
\emph{faithfulness}: how much of the full model's behaviour it reproduces when
everything outside it is ablated, checked with several ablation schemes and metrics.
Our Stage-4 design follows these three points: edges are tested by path patching
into keys and values at the source's fact positions, routes are indexed by position,
and the resulting head mechanism is evaluated by faithfulness on families held out
from every earlier stage.

\subsection{The gap}
\label{sec:rw-gap}

The two lines of work reviewed above rarely meet on the same heads. Descriptive
criteria propose what a head does from its weights or attention; circuit studies
find what heads do by intervening on a task. For semantic induction heads the two
have not been connected: \citet{ren2024semantic} show that some heads stand out on
their relation index and read this as evidence that the heads ``encode'' relations,
but the index is not checked against the model's behaviour. This work examines the
relation between the descriptive SIH criterion and the causal mechanisms that
support the use of in-context relations. Using head-, group- and path-level
interventions in OLMo-2-7B, we characterize the overlap and the gaps between
identification by RI and contribution to behaviour on a relational task; a
complementary analysis across training checkpoints separates the development of
attention routing from the development of the conditional relation index.
